

\documentclass[11pt]{article}

\usepackage[margin=1in]{geometry}
\usepackage{amsmath,amssymb,amsthm,mathtools}
\usepackage{graphicx}
\graphicspath{{figs/}}
\usepackage{booktabs}
\usepackage{array,tabularx}
\usepackage{enumitem}
\usepackage{xurl}
\usepackage{hyperref}
\usepackage{microtype}
\newtheorem{definition}{Definition}
\newtheorem{theorem}{Theorem}
\newtheorem{lemma}{Lemma}
\newtheorem{corollary}{Corollary}
\newtheorem{proposition}{Proposition}
\newtheorem{remark}{Remark}

\newcommand{\Adv}{\mathrm{Adv}}
\newcommand{\TV}{\mathrm{TV}}
\newcommand{\E}{\mathbb{E}}
\newcolumntype{L}[1]{>{\raggedright\arraybackslash}p{#1}}

\title{State-Dependent Anonymity:\\Selection and State Evolution as an Attack Surface}
\author{}
\date{March 22, 2026 (archive v551)}


\begin{document}
\maketitle

\begin{abstract}
Anonymity mechanisms for overlays and anonymous DHTs are increasingly \emph{stateful}:
routing and contact decisions depend on reputations, congestion signals, caches, and operator policies.
State improves performance and security, but it creates a control-plane channel in which tiny per-epoch
selection biases can amplify over time, and in which public state readouts (scores, load reports, cache
summaries) can themselves become witnesses.
We formalize two indistinguishability games---\emph{selection privacy} and \emph{state-evolution privacy}---and give sharp lower bounds
showing how partial observation plus repetition yields near-certain identification.
We then give stability/accounting upper bounds (DP-style) that translate implementable control-plane
constraints into adversary advantage bounds and effective anonymity set sizes.
Finally we isolate a practical mitigation pattern: \emph{state privatization}, in which
selection kernels, state releases, and monitoring are compiled into a single state-anonymity budget.
This note is the theorem/accountant root for the state-anonymity family and is intended to stand alone as its public entrypoint.
The paper stops at that layer: CPPC manifests and alert-publication controls belong to
AnonDHT~State~3, evaluator notarization belongs to Eval~1, and public lineage / release packaging belong to
ABOM/OINL plus Release~A. This paper is designed to be non-redundant with the prior series on stop-time padding,
scheduling/compilers for metadata hiding, and spectral delegation; we cite those works when needed.
\end{abstract}

\paragraph{Compilation note.}
This paper is brief-by-design. A companion addendum collects extended calibration vignettes and worked time-to-identification conversions (State~1A).\cite{series:state1addendum}

\paragraph{Scope.}
We isolate and formalize \emph{State-Dependent Anonymity:} as a reusable theorem/accountant surface in the anonymity / Anonymous-DHT series.
This paper owns the control-plane games, amplification lower bounds, and accountant-to-endpoint conversions for stateful overlays and anonymous DHTs.
It is the public-facing theorem root for that family: once the state claim tuple is declared, the paper can be cited without also carrying the monitoring-manifest or release-lineage layers.
It does not own CPPC manifests or alert-publication policy (AnonDHT~State~3), evaluator notarization / anti-grinding evidence release (Eval~1), or public lineage / publication packaging (ABOM/OINL plus Release~A).\cite{series:state3,series:evalnotary,series:release}

\paragraph{Units.} Budgets are published in bits (use $\log_2$ and $2^{\epsilon}$); results stated in nats can be converted by dividing by $\ln 2$ (see Synthesis~8).\cite{series:notation}

\paragraph{Imports.}
This paper budgets \emph{control-plane state} (scores, caches, load signals) as a witness family.
We treat stop-time padding and schedule compilation as imported from the backbone wiki posts and cite them where needed.
We also treat the endpoint metric interface (odds inflation and identification sizing) as imported; see Anonymity~A and Synthesis~5.\cite{series:oddsinflation,series:endpointbridge}
Trace-interface naming (exposure normal forms, ENF-ID) is imported from Synthesis~3, and threat/repetition windows (TW declarations, TW-ID) from Synthesis~4.\cite{series:traceifaces,series:threatwindows}
Contract-object vocabulary follows Synthesis~0 and the ABOM/OINL note (Anonymity~C), but packaging itself is deferred to those release-facing layers.\cite{series:contractobjects,series:abomoinl,series:release}
For certificate/log mechanics and artifact availability, see the series artifact policy.\cite{series:artifactpolicy}
Notation conventions for keys, state, transcripts, and transcript projections follow the series notation sheet (Synthesis~8).\cite{series:notation}


\paragraph{Exports.}
(i) Two control-plane indistinguishability games (selection privacy and state-evolution privacy) for stateful anonymous overlays and Anonymous-DHT lookups;
(ii) amplification lower bounds showing how small per-epoch selection gaps become near-certain identification under repetition;
(iii) stability/accounting conditions (DP-style) that compile implementable control-plane constraints into odds-inflation / identification budgets; and
(iv) a minimal state-anonymity claim tuple binding \texttt{exposure\_nf\_id}, \texttt{tw\_id}, \texttt{state\_decl\_id}, the chosen accountant, and the resulting endpoint budget/horizon label for later notarization and release packaging.\cite{series:contractobjects,series:abomoinl}
The paper deliberately stops before CPPC manifests, alert-publication policy, receipt issuance, or release-lineage decisions, which belong to AnonDHT~State~3, Eval~1, and Release~A.\cite{series:state3,series:evalnotary,series:release}


\paragraph{Claim fields (non-negotiable).}
All games and bounds in this paper are parameterized by
(i) a declared transcript projection (by \texttt{exposure\_nf\_id}, Synthesis~3) and
(ii) a declared repetition/threat window (by \texttt{tw\_id}, Synthesis~4).\cite{series:traceifaces,series:threatwindows}
For stateful deployments, the public claim also carries a third claim field:
a digest-bound \texttt{state\_decl\_id} committing to a canonical state-surface declaration (what long-lived control-plane state surfaces are in scope and what is excluded).\cite{series:planstateequiv,series:workedexample}
This prevents ``same \texttt{state\_contract\_id}, different state surface'' drift from hiding off-surface.

\paragraph{Validity boundary for a public state-anonymity card.}
A state-anonymity packet is not a blanket anonymity guarantee for deployed traffic.
It covers only the measurable observation channel named by the declared \texttt{exposure\_nf\_id}, \texttt{tw\_id}, and \texttt{state\_decl\_id}.
A release-facing card must therefore say which secret pairs/scenarios are compared; which actions, releases, monitors, cache summaries, scores, retries, and public alerts are in scope; whether the compared scenarios share support; which adaptive kernel changes are accountant-controlled; and which operational knobs are outside the claim.
Padding machines, congestion/path-selection weights, reputation updates, token issuance/redemption surfaces, cache expiration, parser/channel changes, and monitoring dashboards are either part of the declared state surface or excluded from the sentence being certified.
If any of those interfaces changes, the old card is stale unless a successor state card, conditional certificate, or release wrapper explicitly preserves the same claim fields.\cite{series:state3,series:condcert,series:deployedsurfaces,series:planstateequiv,kifer2014pufferfish}

Any later notarized certificate or release receipt should therefore carry \texttt{exposure\_nf\_id}, \texttt{tw\_id}, and (when applicable) \texttt{state\_decl\_id} as explicit line-item claim fields (Synthesis~12), and package the same ids in ABOM/OINL (Anonymity~C).\cite{series:receiptitems,series:abomoinl}
Replay-plan drift is handled analogously via \texttt{plan\_spec\_id}; equivalence/drift rules for both ids are centralized in Synthesis~18.\cite{series:planstateequiv}

\paragraph{Earliest sufficient stop.}
The minimal public object exported by this paper is a declared state-anonymity packet: transcript projection (\texttt{exposure\_nf\_id}), repetition/threat window (\texttt{tw\_id}), state-surface declaration (\texttt{state\_decl\_id} when applicable), game/accountant choice, and the resulting endpoint budget or horizon label.
Once that tuple is fixed, this theorem/accountant note has reached its earliest sufficient public stop. Downstream papers may notarize it (Eval~1), monitor it behind a CPPC/measurement-firewall boundary (AnonDHT~State~3), or package it into lineage-bearing releases (ABOM/OINL + Release~A), but those layers are citation companions rather than prerequisites here. If the live issue is the exact public/on-request packet that should carry one exported state-anonymity card onward, widen only to the maintained worked example; if the live issue is theorem-root-to-review routing across that already-fixed card, widen only to the maintained series-spine bridge.\cite{series:state3,series:evalnotary,series:abomoinl,series:release,series:workedexample,series:seriesspines}

\paragraph{Later-terse-paper citation posture.}
Later terse papers should cite this note only when the live issue is the stateful control-plane theorem/accountant itself: which declared transcript projection, threat window, and state-surface declaration are in scope; which state-anonymity game or stability accountant was used; or which amplification / state-to-endpoint conversion produced the published budget.
If the live issue is instead CPPC manifests or alert-publication controls, stop at AnonDHT~State~3; if it is evaluator notarization, stop at Eval~1; and if it is public claim packaging or lineage, stop at ABOM/OINL and Release~A rather than widening this theorem root into a monitoring or release dossier. If the live issue is the exact public/on-request packet carrying one exported state-anonymity card, widen only to the maintained worked example; if it is theorem-root-to-review routing across that already-fixed card, widen only to the maintained series-spine bridge.\cite{series:state3,series:evalnotary,series:abomoinl,series:release,series:workedexample,series:seriesspines}

\paragraph{State-family branched route.}
The state-anonymity family is easiest to cite as a branched route rather than as one long linear stack.  This note is the first stop and owns the generic stateful theorem/accountant card.  State~1A is only the optional calibration branch when a later paper needs cadence-to-time or horizon conversions for an already-declared State~1 packet.  State~2 is the anonymous-DHT specialization branch when the live issue is instead a committee-contact witness and MUCC-style availability floor for the same deployment slice.  State~3 is the monitoring stop layered above either branch once the question becomes public CPPC fields, alert caps, or measurement-firewall policy, and State~3A is only the optional support packet beneath that monitoring card.  Later terse papers should therefore widen outward by live question, not by habit: theorem root here, optional conversion at State~1A, optional committee-contact specialization at State~2, monitoring at State~3, and optional schema/linter/many-detector support at State~3A.\cite{series:state1addendum,series:state2,series:state3,series:state3addendum}

\paragraph{Tiny branched-route drill.}
Suppose a later short paper needs to say three things at once: the declared state-anonymity packet spends about $0.144$ bits per epoch and $2.31$ bits across the published $16$-epoch horizon; the anonymous-DHT deployment slice also fixed a committee-contact floor of at least $152$ expected committees at $M=256$; and the deployed monitoring surface later published a CPPC with a $6.0$-bit delayed-alert cap.  The first sentence belongs here.  The second belongs to State~2 because it is a DHT-specific committee-contact specialization rather than a generic state-accountant sentence.  The third belongs to State~3 because it is a monitoring-contract sentence rather than a theorem-root sentence.  That is the smallest honest family-level split.\cite{series:state2,series:state3}

\paragraph{A minimal state-anonymity stop card.}
\begin{center}\small
\begin{tabular}{@{}L{3.5cm}L{7.9cm}@{}}
\toprule
\textbf{Live issue} & \textbf{First sufficient owner}\\
\midrule
Which declared \texttt{exposure\_nf\_id} / \texttt{tw\_id} / \texttt{state\_decl\_id}, accountant, and exported budget or horizon row are in force & this note\\
Which alternate horizon table or time-to-identification conversion should be quoted for the same packet & State~1A\\
Which committee-contact floor or abort-safe equalization constraint still governs the same deployment slice & State~2 and Operational~A\\
Which CPPC / alert-publication policy / runtime-conformance receipt attests that slice in deployment & State~3 and Operational~B\\
Which optional schema / linter / many-detector support packet explains that same CPPC card without replacing it & State~3A\\
Which notarized evaluation bundle or lineage-bearing release packet carries the claim onward & Eval~1, Anonymity~C, and Release~A\\
\bottomrule
\end{tabular}
\end{center}
This stop card is the intended fail-closed map for later terse papers: once the live issue has crossed into calibration tables, committee-contact floors, monitoring receipts, evaluator bundles, or release wrappers, keep this theorem/accountant note narrow and cite the first sufficient downstream owner instead of widening State~1 into a general deployment dossier.\cite{series:state1addendum,series:state2,series:operationalA,series:state3,series:operationalB,series:evalnotary,series:abomoinl,series:release}

\paragraph{Artifact policy.}
This archive is source-only; supporting artifacts (plots, scripts, datasets, and tooling) are available on request.\cite{series:artifactpolicy}



\section{Introduction}
Stateful optimizations are now a default feature of anonymous overlays and anonymous DHT-style systems.
Selection and routing decisions depend on \emph{evolving} signals such as reputation, congestion, and caches;
operators also run monitors and publish summaries.
This paper treats that statefulness as a first-class anonymity channel: the \emph{control plane} becomes a stochastic process whose visible actions and releases can leak secrets even when payloads are protected.

\paragraph{Control-plane loop.}
Figure~\ref{fig:control-plane-loop} highlights the structural pattern.
A secret (destination, user profile, workload class) influences hidden state, state influences selections, and selections feed back into state.
An adversary need not observe payloads: partial visibility of selections, state readouts, or induced probes can suffice.

\begin{figure}[t]
  \centering
  \fbox{\parbox{0.92\linewidth}{\centering\textbf{Figure omitted in source-only archive.}\\(Plot/diagram and scripts available on request.)}}
  \caption{State-dependent anonymity as a control-plane feedback loop.
  Small per-epoch selection biases (from state) can amplify over time; public state readouts can be witnesses.
  }
  \label{fig:control-plane-loop}
\end{figure}

\paragraph{Scope and non-redundancy.}
We focus on \emph{state-dependent} leakage.
We treat stop-time padding and schedule equalization as imported mechanisms from the companion series~\cite{series:stoptimeaddendum,series:schedcompiler};
we use spectral delegation results when we need mixing identities~\cite{series:spectral}.
For DHT-specific contact-set floors (MUCC), see AnonDHT~State~2.\cite{series:state2}
For deployment-side monitoring, alarms, and CPPC manifests, see AnonDHT~State~3.\cite{series:state3}

Deployed control-plane configuration surfaces (e.g., delegated-routing endpoint sets, AutoConf resolution, and OHTTP/chunking feature drift) are treated separately in Synthesis~30 and should be pinned via \texttt{state\_decl\_id} / separation manifests rather than left implicit.\cite{series:deployedsurfaces}

\paragraph{Contributions.}
\begin{itemize}[leftmargin=*]
\item \textbf{Games:} we define \emph{selection privacy} and \emph{state-evolution privacy} for control-plane processes.
\item \textbf{Lower bounds:} witness-style distinguishers and repetition/amplification theorems that quantify how tiny per-epoch gaps become near-certain identification.
\item \textbf{Upper bounds:} stability/accounting conditions (DP-style) that translate implementable control-plane constraints into advantage bounds and effective anonymity set guarantees.
\item \textbf{Mitigations:} a practical \emph{state-privatization} pattern and a refresh/expiration perspective for making leakage time-limited; CPPC manifests, alert publication, and runtime monitoring stay in the AnonDHT~State~3 companion.\cite{series:state3}
\end{itemize}


\section{Model and anonymity games}

\paragraph{Parameterization by interface and window.}
Let \(\mathsf{ENF}\) denote the exposure normal form that names the transcript projection revealed to the adversary, and let \(\mathsf{TW}\) denote the repetition/threat window.
We write their claim fields as \texttt{exposure\_nf\_id} and \texttt{tw\_id} (full digests in ABOM).
When we say ``\(n\) epochs'' or ``a horizon of \(n\) lookups'', that \(n\) is the window field of \(\mathsf{TW}\), and the adversary's view is always the projection named by \(\mathsf{ENF}\).

\label{sec:games}

\paragraph{Protocol as a process.}
Fix a secret $\sigma \in \{0,1\}$ (or more generally $\sigma\in[K]$).
A stateful protocol induces a stochastic process $(S_t,A_t,R_t)_{t\ge 1}$ where:
(i) $S_t$ is internal state (scores, queues, cache contents, committee tables),
(ii) $A_t$ is the control-plane action (selected relay/committee/contact-set, schedule choice), and
(iii) $R_t$ is any \emph{released} state (published weights, beacons, alerts).
The adversary observes $O_t = \phi(A_t,R_t)$ for some observation map $\phi$ capturing partial compromise, sampling, or probing.

\paragraph{Selection privacy.}
Selection privacy asks what an observer learns from the distribution of chosen control-plane actions.

\begin{definition}[Selection-privacy experiment]
Let $\mathsf{View}_\sigma$ denote the distribution of the adversary transcript $O_{1:T}$ when the secret is $\sigma$.
The selection advantage is
\[
\Adv^{\mathsf{sel}}_T \;:=\; \TV\big(\mathsf{View}_0,\mathsf{View}_1\big),
\]
where $\TV(P,Q)=\sup_E |P(E)-Q(E)|$.
\end{definition}

This definition is agnostic to implementation: it applies whether actions are relay choices, DHT committee contacts, or schedule selections.

\begin{remark}[MUCC as a marginal instantiation (DHT contact sets)]
If actions encode a (possibly multi-round) committee contact set in an anonymous DHT, then a common design goal is that the destination key does not affect which committees are contacted.
Paper~2 studies a weak, marginal-only version of this objective (MUCC): each individual committee is contacted with key-independent probability~\cite{series:state2}.
MUCC does \emph{not} imply selection privacy in the above sense (the joint distribution of the contacted set may still depend on the key), but it is enough to force strong availability/contact floors.
\end{remark}

\paragraph{State-evolution privacy.}
State-evolution privacy isolates leakage from public state readouts that drive future choices.

\begin{definition}[State-evolution-privacy experiment]
Let $\mathsf{Rel}_\sigma$ be the transcript distribution when the observer sees only releases $R_{1:T}$ (and any derived metadata) under secret $\sigma$.
Define
\[
\Adv^{\mathsf{state}}_T \;:=\; \TV\big(\mathsf{Rel}_0,\mathsf{Rel}_1\big).
\]
\end{definition}

\paragraph{Coupling view: observation as a channel.}
Partial observability is an additional channel applied to the underlying action/state process.
This can \emph{amplify} or \emph{contract} distinguishability depending on what is exposed.
A useful (and implementable) abstraction is to treat observation as subsampling or coarsening, which often yields privacy amplification in DP analyses.
We use this later as an upper-bounding tool.

\paragraph{Beyond binary secrets.}
For $\sigma\in[K]$, the natural goal is to upper bound the maximum a-posteriori (MAP) identification probability.
In Section~\ref{sec:upper}, we relate stability to odds inflation and to an effective anonymity set $K_{\mathrm{eff}}$.


\section{Witness lower bounds and amplification}
\label{sec:lower}

A core lesson of state-dependent anonymity is that \emph{any} object whose selection probability differs across secrets is a witness.
Repeated exposure turns tiny gaps into near-certain identification.

\subsection{Single-witness distinguishers}
\label{sec:witness}
Let $W$ be a per-epoch event in the observer view (e.g., ``committee $j$ was contacted'', ``relay $r$ was selected'', ``cache hit observed'').
Let $p_\sigma := \Pr[W\mid \sigma]$.

\begin{theorem}[Witness bound]
\label{thm:sel-witness}
For any horizon $T$, the selection advantage satisfies
\[
\Adv^{\mathsf{sel}}_T \;\ge\; \TV\big(\mathrm{Bin}(T,p_0),\mathrm{Bin}(T,p_1)\big),
\]
whenever the observer can extract (possibly via a measurable function of $O_{1:T}$) an i.i.d. witness indicator sequence $(\mathbf{1}\{W_t\})_{t=1}^T$ with per-epoch probabilities $p_0,p_1$.
\end{theorem}

\begin{proof}
Let $g$ map the full transcript $O_{1:T}$ to the witness count $N:=\sum_{t=1}^T \mathbf{1}\{W_t\}$.
By assumption, under secret $\sigma$ the induced distribution of $N$ is $\mathrm{Bin}(T,p_\sigma)$.
Total variation distance cannot increase under post-processing (measurable maps), so
\[\TV\big(\mathsf{View}_0,\mathsf{View}_1\big)\;\ge\;\TV\big(g(\mathsf{View}_0),g(\mathsf{View}_1)\big)\;=\;\TV\big(\mathrm{Bin}(T,p_0),\mathrm{Bin}(T,p_1)\big).\]
\end{proof}

The i.i.d. condition holds in many ``static secret'' analyses (fixed destination / profile) and is a conservative approximation in slowly drifting regimes.

\subsection{Repetition and sample complexity}
\label{sec:amplification}
Write $\Delta := |p_1-p_0|$.
Testing two Bernoulli means with a sequential probability ratio test (SPRT) yields sharp amplification rates~\cite{wald_sprt_1945,howard2021confseq}.

\begin{theorem}[Repeated-witness amplification]
\label{thm:repeat-witness}
Assume witness indicators $X_1,\dots,X_T$ are i.i.d. Bernoulli$(p_\sigma)$ under secret $\sigma\in\{0,1\}$.
There exists a test with total error at most $\eta$ whenever
\[
T \;\gtrsim\; \frac{\log_2(1/\eta)}{\Delta^2}.
\]
Conversely, for $\Delta$ small and $p_\sigma$ bounded away from $0,1$, any test needs $T=\Omega(\log_2(1/\eta)/\Delta^2)$ samples.
\end{theorem}

Figure~\ref{fig:eps-to-samples} gives a ``numbers-first'' view of this scaling.

\begin{figure}[t]
  \centering
  \fbox{\parbox{0.92\linewidth}{\centering\textbf{Figure omitted in source-only archive.}\\(Plot/diagram and scripts available on request.)}}
  \caption{Illustrative sample complexity scaling: tiny per-epoch witness gaps require long horizons, but eventually yield near-certain identification.
  }
  \label{fig:eps-to-samples}
\end{figure}

\paragraph{Calibration snapshot (informal).}
Table~\ref{tab:calibration} gives a compact numbers-first view of the $1/\Delta^2$ amplification scaling in Theorem~\ref{thm:repeat-witness}.
The table is illustrative (independence/nonstationarity ignored) and exists to make the quadratic dependence operationally visible.

\begin{table}[h]
\centering
\footnotesize
\setlength{\tabcolsep}{4pt}
\begin{tabularx}{\linewidth}{@{}r r r >{\raggedright\arraybackslash}X@{}}
\toprule
Per-epoch gap $\Delta$ & Target total error $\eta$ & Rule-of-thumb $T$ & Interpretation \\
\midrule
0.01 & 0.05 & $\approx 6\times 10^4$ & ``tiny bias'' needs $\sim$weeks at minute cadence (or days at 10-second cadence). \\
0.02 & 0.05 & $\approx 1.5\times 10^4$ & $\sim$10 days at minute cadence (or $\sim$2 months at 5-minute cadence). \\
0.05 & 0.05 & $\approx 2.4\times 10^3$ & $\sim$40 minutes at 1-second cadence (or $\sim$7 hours at 10-second cadence). \\
0.10 & 0.05 & $\approx 6\times 10^2$ & $\sim$10 minutes at 1-second cadence (or $\sim$1.2 hours at 10-second cadence). \\
\bottomrule
\end{tabularx}
\normalsize
\caption{Back-of-the-envelope calibration for Theorem~\ref{thm:repeat-witness} using $T\approx 2\log_2(1/\eta)/\Delta^2$.
These are illustrative and ignore dependence and nonstationarity; they serve only to show the strong $1/\Delta^2$ scaling.
}
\label{tab:calibration}
\end{table}


\subsection{Adaptive witnesses and many streams}
Real deployments generate many potential witnesses (keys, committees, relays, cache lines, alarms).
Even when any \emph{single} stream is weak, scanning or adaptively choosing witnesses can drive identification.

\begin{corollary}[Searching $N$ witness streams adds a $\log_2 N$ factor]
\label{cor:many-witness-search}
Suppose the observer can compute $N$ candidate witness streams
$\{X^{(i)}_t\}_{t\ge 1}$ for $i\in[N]$, where each $X^{(i)}_t\in\{0,1\}$ and (under a fixed secret $\sigma\in\{0,1\}$) each stream is i.i.d.\ Bernoulli$(p_{\sigma,i})$.
Let $\Delta_\star := \max_i |p_{1,i}-p_{0,i}|$.
Then there exists an anytime-valid sequential test that (i) adaptively chooses a stopping time $\tau$,
(ii) has total error at most $\eta$, and (iii) achieves the scaling
\[
\tau \;=\; O\!\left(\frac{\log_2(N/\eta)}{\Delta_\star^2}\right)
\]
up to constant factors (in regimes where $p_{\sigma,i}$ are bounded away from $0,1$).
\end{corollary}

\begin{proof}[Proof sketch]
Run a likelihood-ratio e-process (SPRT-style) on each candidate stream in parallel and stop when any crosses a calibrated threshold.
Ville's inequality gives time-uniform control under optional stopping, and the extra $\log_2 N$ factor comes from the multiple-testing calibration over $N$ streams; see the SAVI/e-value framework~\cite{ramdas2024evalues,wang2025sebH}.
\end{proof}

We avoid duplicating the full many-testing development in the companion series paper~\cite{series:manytestinglowerbounds}.
Instead, Paper~3 in this archive uses e-values/anytime-valid monitoring as an \emph{accounting interface} for many public streams (alerts, beacons, summaries).


\section{Stability/accounting upper bounds}
\label{sec:upper}

Lower bounds say what can go wrong.
Upper bounds should be implementable: they must correspond to constraints engineers can enforce on selection kernels, releases, and monitors.
A natural language is \emph{stability}: small changes in the secret should not strongly tilt control-plane probabilities.

\subsection{DP-style stability for control planes}
\label{sec:dpstyle}

\begin{definition}[Per-epoch max-divergence stability, bit units]
A per-epoch observation distribution $\mathsf{Obs}_\sigma$ satisfies $(\varepsilon,\delta)$ stability if for all measurable $E$,
\[
\Pr[\mathsf{Obs}_0\in E] \le 2^{\varepsilon}\Pr[\mathsf{Obs}_1\in E]+\delta
\quad\text{and}\quad
\Pr[\mathsf{Obs}_1\in E] \le 2^{\varepsilon}\Pr[\mathsf{Obs}_0\in E]+\delta.
\]
\end{definition}

This is the standard (approximate) differential privacy inequality~\cite{dworkroth2014foundations} applied to control-plane views, written in the bit-unit convention used by this archive. A natural-log likelihood-ratio cap $\exp(\varepsilon_{\mathrm{nat}})$ must therefore be converted to $\varepsilon_{\mathrm{bits}}=\varepsilon_{\mathrm{nat}}/\ln 2$ before it is inserted into this definition.
It immediately yields a TV bound and composes over time.

\begin{proposition}[Stability implies TV and odds bounds]
\label{prop:dp-to-tv-odds}
If $\mathsf{Obs}_\sigma$ satisfy $(\varepsilon,\delta)$ stability, then
\[
\TV\big(\mathsf{Obs}_0,\mathsf{Obs}_1\big) \le \min\{1,\,2^{\varepsilon}-1+\delta\}.
\]
If $\delta=0$ and the distributions admit densities with respect to a common measure, then for any realized observation $o$,
\[
2^{-\varepsilon} \le \frac{d\mathsf{Obs}_0}{d\mathsf{Obs}_1}(o) \le 2^{\varepsilon},
\]
so posterior odds between $\sigma=0$ and $\sigma=1$ inflate by at most $2^{\varepsilon}$ in either direction.
Under basic composition over epochs with per-epoch $(\varepsilon_t,\delta_t)$, the transcript satisfies $(\sum_t \varepsilon_t,\sum_t \delta_t)$, and the same bounds apply with $\varepsilon=\sum_t\varepsilon_t$ and $\delta=\sum_t\delta_t$.
\end{proposition}

\begin{proof}[Proof sketch]
Let $E$ attain the TV supremum so that $\TV(P,Q)=P(E)-Q(E)$ for $P=\mathsf{Obs}_0,Q=\mathsf{Obs}_1$.
Apply the stability inequality to $E$ and to its complement and rearrange; for example,
$P(E)-Q(E)\le (2^{\varepsilon}-1)Q(E)+\delta\le 2^{\varepsilon}-1+\delta$, and the symmetric inequality handles the other sign.
The odds statement for $\delta=0$ is the pointwise likelihood-ratio form of max-divergence stability, and basic composition follows by multiplying likelihood-ratio envelopes (or summing the usual approximate-DP losses) over the declared epochs.
\end{proof}

\subsection{Stopping-safe drift audits (optional stopping)}
Control planes are monitored in real time: operators run drift detectors on selections, scores, and release channels, and trigger mitigation when the detector fires.
But the trigger time is \emph{adaptive} (a stopping time), and naive ``check every epoch'' thresholds can severely understate false-alarm rates.
A clean and receipt-friendly fix is to publish drift evidence as a \emph{time-uniform} bound derived from a nonnegative supermartingale (an \emph{e-value}).\cite{howard2018cher,wang2020ebh}

\begin{lemma}[Ville inequality / time-uniform validity]
Let $(E_t)_{t\ge 0}$ be a nonnegative supermartingale with $E_0=1$ under a null model (e.g., ``no drift''), and let $\tau$ be any (possibly data-dependent) stopping time.
Then for any $\alpha\in(0,1)$,
\[
\Pr\Big[\sup_{t\ge 0} E_t \ge 1/\alpha\Big] \le \alpha,
\quad\text{and hence}\quad
\Pr\big[E_{\tau}\ge 1/\alpha\big]\le \alpha.
\]
\end{lemma}

\begin{remark}[Receipt hook]
A drift monitor can therefore publish: (i) the detector family (as part of a state declaration), (ii) the chosen $\alpha$, and (iii) a digest-bound evidence object containing the realized $E_t$ trace up to $\tau$.
Because the guarantee is time-uniform, auditors do not need to know how often the monitor checked or how the trigger time was chosen.
This is orthogonal to the privacy bound itself: if a triggered mitigation changes the control-plane kernel, then the updated receipt should roll the relevant envelope ids (new \texttt{state\_decl\_id} and/or new knobs), and any adaptive accounting should be expressed via conditional certificates (Synthesis~21).\cite{series:condcert,series:planstateequiv}
\end{remark}


\subsection{Concrete kernel: softmax stability from bounded score sensitivity}
Many control planes implement selection via a scored softmax (or, equivalently, the exponential mechanism~\cite{mcsherry2007mechanism}).
Let scores $x_i(\sigma)\in\mathbb{R}$ be computed for each option $i\in[M]$, and select $I\in[M]$ with
\[
\Pr[I=i\mid \sigma] \;=\; \frac{\exp(x_i(\sigma)/\tau)}{\sum_{j=1}^M \exp(x_j(\sigma)/\tau)},
\]
where $\tau>0$ is a temperature.

\begin{lemma}[Softmax max-divergence bound, unit-normalized]
\label{lem:softmax-stability}
Assume $\|x(\sigma{=}0)-x(\sigma{=}1)\|_\infty \le \Delta$.
Then the selection distribution satisfies $(\varepsilon_{\mathrm{bits}},0)$ stability in the bit-unit convention used in Proposition~\ref{prop:dp-to-tv-odds}, where
\[
\varepsilon_{\mathrm{nat}} \;=\; \frac{2\Delta}{\tau},
\qquad
\varepsilon_{\mathrm{bits}} \;=\; \frac{2\Delta}{\tau\ln 2}.
\]
\end{lemma}

\begin{proof}
For each $i$, the numerator ratio is bounded by $\exp(\Delta/\tau)$, while the normalizer ratio is bounded by $\exp(\Delta/\tau)$ in the opposite direction:
\[
\frac{\sum_j \exp(x_j(1)/\tau)}{\sum_j \exp(x_j(0)/\tau)} \le \exp(\Delta/\tau).
\]
Thus $\Pr[I=i\mid 0]/\Pr[I=i\mid 1]\le \exp(2\Delta/\tau)$ for every $i$. This is a nat-unit likelihood-ratio cap $\varepsilon_{\mathrm{nat}}=2\Delta/\tau$; because Proposition~\ref{prop:dp-to-tv-odds} uses $2^{\varepsilon}$ odds notation, the release-card value is $\varepsilon_{\mathrm{bits}}=\varepsilon_{\mathrm{nat}}/\ln 2$.
\end{proof}

Lemma~\ref{lem:softmax-stability} turns an engineering control (score sensitivity $\Delta$ and temperature $\tau$) into an explicit per-epoch stability budget.
In practice, $\Delta$ can be enforced by clipping or smoothing state-dependent scores before selection, and by limiting how much any single secret-dependent update can move a score.

\paragraph{Partial observation as privacy amplification.}
When the observer sees only a subsample or coarsening of actions, stability can improve.
We treat this using privacy amplification by subsampling~\cite{balle2018subsampling} as a conservative upper bound.

\subsection{Composition with adaptive parameters}
In stateful systems, selection parameters and monitoring intensity may be chosen adaptively.
Stability budgets must therefore be tracked with an \emph{online accountant}.
Privacy odometers/filters provide a canonical interface for this pay-as-you-go regime~\cite{rogers2016odometers,whitehouse2023fullyadaptive}.
When multiple mechanisms interleave (e.g., routing + monitoring + retries), interactive concurrent composition is relevant~\cite{vadhanwang2021concurrentdp,zhang2023concurrentinteractive}.
We defer the deployment-side discussion (budget beacons, alarm release, and practical filter choices) to Paper~3~\cite{series:state3}.
\subsection{Effective anonymity sets from odds inflation}
For $K$ profiles, a compact operational summary is an \emph{effective anonymity set}.
A simple and widely useful interpretation is via worst-case odds inflation.
If stability guarantees that the likelihood ratio between any two profiles is at most $\rho$ (over the observed horizon), then MAP identification is at most $\rho/K$ under a uniform prior, suggesting
$K_{\mathrm{eff}}\approx K/\rho$.
Figure~\ref{fig:keff-vs-eps} shows a typical calibration between an $\varepsilon$-style budget and $K_{\mathrm{eff}}$.

\begin{figure}[t]
  \centering
  \fbox{\parbox{0.92\linewidth}{\centering\textbf{Figure omitted in source-only archive.}\\(Plot/diagram and scripts available on request.)}}
  \caption{Example calibration between $\varepsilon$-style stability and effective anonymity size $K_{\mathrm{eff}}$.}
  \label{fig:keff-vs-eps}
\end{figure}

\subsection{Expiration and refresh}
\label{sec:expiration}
If a system mixes and periodically refreshes state (or applies gradual privacy expiration), then leakage can become time-limited.
We treat this as ``expiration-aware'' accounting and cite discounted DP and gradual expiration mechanisms~\cite{Farokhi2020DiscountedDP,AnderssonEtAl2024GradualExpiration}.

\subsection{Contractive state updates justify expiration-style models}
\label{sec:contraction}
Many practical control-plane states are implemented with exponential forgetting (EWMA reputations, cache decay, moving averages).
This yields a clean contraction bound that turns an ``infinite horizon'' into an \emph{effective} finite one.

\begin{lemma}[EWMA contraction]
\label{lem:ewma-contraction}
Let a scalar state $s_t\in\mathbb{R}$ evolve as
\[
s_{t+1}=(1-\lambda)s_t+\lambda u_t,\qquad \lambda\in(0,1],
\]
where $u_t$ is an arbitrary input sequence.
For any two executions with the same inputs but different initial states $s_0,s_0'$, the difference contracts geometrically:
\[
|s_t-s_t'|\le (1-\lambda)^t\,|s_0-s_0'|.
\]
More generally, if the inputs differ by at most $|u_t-u_t'|\le \Delta_u$ for all $t$, then
\[
|s_t-s_t'|\le (1-\lambda)^t|s_0-s_0'|+\Delta_u\sum_{i=0}^{t-1}(1-\lambda)^i\le (1-\lambda)^t|s_0-s_0'|+\Delta_u/\lambda.
\]
\end{lemma}

\begin{proof}
Unroll the recursion and sum the resulting geometric series.
\end{proof}

\begin{remark}[Vector-state variant]
If the state is vector-valued $s_t\in\mathbb{R}^p$ and evolves via the same EWMA recursion $s_{t+1}=(1-\lambda)s_t+\lambda u_t$, then the same geometric contraction holds in any norm:
$\|s_t-s_t'\|\le (1-\lambda)^t\|s_0-s_0'\|$ (plus the corresponding input-difference term).
If a selection kernel is Lipschitz in state in the sense that $\TV(K(s),K(s'))\le L\|s-s'\|$ (or via an SDPI/Dobrushin coefficient for the induced channel), then contraction yields an explicit ``effective memory'' bound that is directly usable in expiration-aware accounting.
\end{remark}


If a selection kernel $K(\cdot)$ is Lipschitz in state (e.g., $\TV(K(s),K(s'))\le L|s-s'|$ or a comparable max-divergence notion), then Lemma~\ref{lem:ewma-contraction} implies that state differences from ancient epochs are exponentially downweighted.
Operationally, $1/\lambda$ acts as an ``effective memory'' in epochs, aligning with discounted/gradual-expiration accounting.


\section{Mitigation pattern: state privatization}
\label{sec:mitigations}

A practical lesson is that control-plane choices, public releases, and monitoring must be designed \emph{together}.
Treating them independently tends to create accidental witnesses.

\paragraph{State privatization compiler.}
We advocate a compiler viewpoint:
\emph{(i)} define a target selection kernel family, \emph{(ii)} define which state is released and at what cadence,
\emph{(iii)} define monitoring/alerts, and \emph{(iv)} attach a single budget/accountant to the whole control plane.
Paper~3 in this archive develops a deployment pattern for this idea (measurement firewall + CPPC manifests).

\paragraph{Refresh to make leakage expire.}
If the state process mixes and occasionally refreshes, then older secrets have diminishing influence on future actions.
Figure~\ref{fig:forgetting} illustrates an ``identification half-life'' viewpoint.

\begin{figure}[t]
  \centering
  \fbox{\parbox{0.92\linewidth}{\centering\textbf{Figure omitted in source-only archive.}\\(Plot/diagram and scripts available on request.)}}
  \caption{Illustrative decay/expiration curve for stateful leakage.}
  \label{fig:forgetting}
\end{figure}

\paragraph{Caution: inhomogeneous mixing can leak.}
Not all mixing is benign: non-homogeneous dynamics can create spectral leakage modes.
We rely on the series' spectral delegation paper for detailed identities~\cite{series:spectral};
Figure~\ref{fig:inhom} is a warning example.

\begin{figure}[t]
  \centering
  \fbox{\parbox{0.92\linewidth}{\centering\textbf{Figure omitted in source-only archive.}\\(Plot/diagram and scripts available on request.)}}
  \caption{Example of inhomogeneous spectral leakage.}
  \label{fig:inhom}
\end{figure}

\paragraph{Checklist.}
In practice, the following questions catch most state-dependent failures:
(i) what state influences selection, and who can observe it;
(ii) which public releases can become probes/oracles;
(iii) how are retries/abort patterns equalized (see stop-time series work~\cite{series:stoptimeaddendum});
(iv) which monitors/alerts exist, and are they closed-view;
(v) what refresh mechanisms bound the influence horizon.


\section{Related work}
Statefulness pervades modern anonymity.
Tor's design emphasizes practical deployability~\cite{dingledine2004tor} and has a long line of work on congestion/selection.
Recent bandwidth-measurement manipulation attacks (e.g., TorMult and MirageFlow) illustrate that published selection weights can be steered, turning the control plane into an oracle for downstream traffic analysis~\cite{Sendner2023TorMult,Sendner2024MirageFlow}.
DPSelect is an early proposal to apply max-divergence accounting directly to Tor guard selection~\cite{Hanley2019DPSelect}.

In incentivized mixnets, routing and rewards are coupled to stake/performance; reputation becomes a direct attack surface.
Recent work analyzes and attacks Nym's reputation system (Cao and Green)~\cite{cao2026nymrep}. The PETS~2026 paper on website fingerprinting on Nym released artifacts/datasets on Zenodo~\cite{jolles2025wf4nym_zenodo}.
Flow-correlation and attribution attacks continue to improve (including deep learning)~\cite{nasr2018deepcorr,oh2022deepcoffea,singh2025revealnet}.
For mixnets specifically, MixMatch demonstrates flow matching at the mixnet level~\cite{oldenburg2024mixmatch}.
For Tor onion services, sliding subset-sum correlation is a recent example of long-horizon correlation under partial observability~\cite{lopes2024sumo}.

Our upper bounds borrow heavily from differential privacy as a stability language~\cite{dworkroth2014foundations}.
We rely on odometers/filters for adaptive parameter regimes~\cite{rogers2016odometers} and on interactive/concurrent composition results when multiple deployed mechanisms interleave~\cite{vadhanwang2021concurrentdp,zhang2023concurrentinteractive}.
We use privacy amplification by subsampling as a conservative model of partial observation~\cite{balle2018subsampling}.
Repeated aggregate releases over time are treated in the continual-observation DP literature~\cite{dwork2010continual,chan2010continual}.
Always-valid monitoring and e-values provide a clean technical interface between detection and accounting~\cite{howard2021confseq,ramdas2024evalues}. We develop the deployment-side implications (including multiple-testing and public-alert leakage) in Paper~3.

For explicitly stateful secrets and correlated background knowledge, Pufferfish provides a general Bayesian framework that treats (i) the protected secrets, (ii) discriminative secret pairs, and (iii) data evolution scenarios as first-class inputs~\cite{kifer2014pufferfish}.
While we keep the present paper in a lightweight max-divergence/DP style, Pufferfish-style instantiations are a natural way to formalize state-evolution privacy when the operator and observer share an explicit model of state dynamics.

Caching as a privacy side channel has a long history (timing/oracle effects) and appears in anonymity, content routing, and DHTs.
For anonymous DHTs, the control-plane state includes routing tables, bucket replacement/refresh, and provider-record lifetimes.
Paper~2 of this split archive treats DHT-specific anonymous lookup mechanisms, contact-set witness floors (MUCC), and deployed-scale context (including recent IPFS evidence).




\paragraph{A minimal public state-anonymity card.}
Table~\ref{tab:minimal-state-anonymity-card} is the smallest public packet later terse papers should need when the live question is simply \emph{which stateful control-plane anonymity contract was declared for this deployment slice?}

\begin{table}[t]
\centering
\small
\begin{tabular}{@{}L{0.33\linewidth}L{0.57\linewidth}@{}}
\toprule
\textbf{Public row} & \textbf{Pinned value}\\
\midrule
Declared observation / window & \texttt{exposure\_nf\_id = enf-committee-contact-bucket-v2}; \texttt{tw\_id = tw-24h-200-lookups} \\
Declared state surface & \texttt{state\_decl\_id = state-routing-scores-cache-v1} covering routing scores and cache summaries, excluding payload contents \\
Chosen accountant & softmax score-sensitivity accountant from Lemma~\ref{lem:softmax-stability} with $\|x(0)-x(1)\|_\infty \le \Delta = 0.01$ and temperature $\tau = 0.20$ \\
Per-epoch stability export & $\varepsilon_{\mathrm{nat}} = 2\Delta/\tau = 0.10$ nats, i.e. $\varepsilon_{\mathrm{bits}} = 0.10/\ln 2 \approx 0.144$ bits \\
Horizon row & under $16$ epochs with no reset, Proposition~\ref{prop:dp-to-tv-odds} gives total $\varepsilon_{\mathrm{bits}} \approx 16\cdot 0.144 = 2.31$ bits by basic composition \\
Public interpretation & on the declared state surface, posterior odds between the two compared state scenarios inflate by at most $2^{0.144} \approx 1.105$ per epoch under the stated sensitivity cap \\
Downstream owner & AnonDHT~State~3 / Eval~1 / ABOM--OINL / Release~A for CPPC manifests, notarized evidence, and lineage-tagged public receipts \\
\bottomrule
\end{tabular}
\caption{A minimal public state-anonymity card. This is the non-decorative public answer later terse papers can quote directly when the live issue is the declared stateful control-plane accountant rather than CPPC monitoring, evaluator notarization, or release lineage.}
\label{tab:minimal-state-anonymity-card}
\end{table}

This card is intentionally narrow: it pins the declared observation/window ids, the state surface, one concrete stability accountant, and the resulting per-epoch / horizon budget rows, but it does not absorb CPPC manifests, monitoring evidence packets, or release receipts. Exact public/on-request packet wiring for one exported state-anonymity card is intentionally delegated to the maintained worked example, and theorem-root-to-review routing across that already-fixed card is delegated to the maintained series-spine bridge. Those remain owned by the downstream companions already cited above.\cite{series:state3,series:evalnotary,series:abomoinl,series:release,series:workedexample,series:seriesspines}

\paragraph{One tiny accountant drill.}
With the card values above, Lemma~\ref{lem:softmax-stability} gives
\[
\varepsilon_{\mathrm{nat}} = \frac{2\Delta}{\tau} = \frac{2\cdot 0.01}{0.20} = 0.10.
\]
Converting to bits yields
\[
\varepsilon_{\mathrm{bits}} = \frac{0.10}{\ln 2} \approx 0.144.
\]
If the same cap applies over $16$ epochs with no reset, Proposition~\ref{prop:dp-to-tv-odds} gives the basic-composition horizon budget
\[
16\cdot 0.144 \approx 2.31\ \text{bits}.
\]
A later terse paper can quote the card and this arithmetic without reopening the lower-bound calibration tables, CPPC manifests, or release packaging layers.

\begin{proposition}[State-card unit conversion guard]
\label{prop:state-card-unit-guard}
The card in Table~\ref{tab:minimal-state-anonymity-card} is release-valid only when the natural-log softmax ratio bound from Lemma~\ref{lem:softmax-stability} is converted into the bit-unit stability parameter used by Proposition~\ref{prop:dp-to-tv-odds}.  Thus the exported row is $\varepsilon_{\mathrm{bits}}=(2\Delta/\tau)/\ln 2$, not the raw nat value $2\Delta/\tau$.  A receipt that copies $0.10$ as a bit exponent, composes $16\cdot 0.10$ bits, or quotes odds $2^{0.10}$ for this card is not source-bound to this theorem; the safe exported values are $0.144$ bits per epoch, $2.31$ bits over the declared $16$ epochs, and per-epoch odds inflation about $1.105$.
\end{proposition}

\paragraph{One tiny worked-route drill.}
Keep the same declared ids and horizon row from Table~\ref{tab:minimal-state-anonymity-card}. Then a later terse paper may honestly say exactly one sentence:
\begin{quote}\small\emph{Under\\
\nolinkurl{exposure_nf_id = enf-committee-contact-bucket-v2},\\
\nolinkurl{tw_id = tw-24h-200-lookups}, and\\
\nolinkurl{state_decl_id = state-routing-scores-cache-v1},\\
the published softmax state-anonymity packet spends about $0.144$ bits per epoch and $2.31$ bits across the declared $16$-epoch horizon.}\end{quote}
That sentence is honest because the maintained worked packet already fixes which public card carries those ids and rows, while theorem-root-to-review routing across that card is maintained by the series-spine bridge rather than reconstructed locally here.\cite{series:workedexample,series:seriesspines}

\paragraph{One tiny paired-route drill.}
Keep the same declared ids and horizon row from Table~\ref{tab:minimal-state-anonymity-card}, and suppose the separately cited State~1A calibration packet fixes \texttt{cadence = 1 witness / hour}, \texttt{eta = 0.05}, and the vignette row \texttt{Delta = 0.02 -> T \textasciitilde 1.5e4 epochs -> 1.7 years @ 1 hour}.
Then a later terse Anonymous-DHT note may honestly pair exactly two short sentences: \emph{``Under \nolinkurl{exposure_nf_id = enf-committee-contact-bucket-v2}, \nolinkurl{tw_id = tw-24h-200-lookups}, and \nolinkurl{state_decl_id = state-routing-scores-cache-v1}, the published state-anonymity packet spends about $0.144$ bits per epoch and $2.31$ bits across the declared $16$-epoch horizon. Under the separately cited State~1A calibration packet, a witness gap $\Delta=0.02$ at $\eta=0.05$ corresponds to about $1.5\times 10^{4}$ epochs, i.e. about $1.7$ years at one witness per hour.''}
That paired route stays non-redundant only because the theorem/accountant claim remains owned here while the cadence-to-time conversion remains owned by State~1A; neither sentence silently replaces the other.\cite{series:state1addendum}

\paragraph{One tiny fail-closed cutover drill.}
Reuse the same card, ids, and accountant. After $12$ epochs on the same declared state surface, the basic-composition spend is only
\[
12\cdot 0.144 \approx 1.73\ \text{bits},
\]
so the old card still licenses the statement ``this deployment slice stayed within the published $16$-epoch state-anonymity packet so far.''
But if the operator quietly extends the same slice to $24$ epochs with no reset, the old packet no longer licenses a $24$-epoch claim: the same arithmetic gives
\[
24\cdot 0.144 \approx 3.46\ \text{bits},
\]
which lies outside the published horizon row. The correct terse-paper move is therefore fail closed under the old $16$-epoch card until a new packet is issued.
Likewise, if epochs $13$--$16$ newly publish a cache-eviction digest or another control-plane summary that was not covered by \texttt{state\_decl\_id}, the old card cannot be stretched by arithmetic alone: the state surface changed, so the old packet is no longer the quotable public object even before the horizon is exhausted.

\paragraph{One tiny same-card / refreshed-wrapper drill.}
Keep the same declared ids, the same \texttt{state\_decl\_id}, the same accountant row, and the same $16$-epoch horizon row, but let the operator issue a successor ABOM / Release~A packet only because the citation head, signature envelope, or lineage metadata changed.
Then the theorem/accountant sentence exported above is still the same public state-anonymity claim: a later terse paper may keep citing State~1 for the row ``about $0.144$ bits per epoch and $2.31$ bits across the declared $16$-epoch horizon,'' and widen only to Anonymity~C or Release~A for the refreshed wrapper.
By contrast, if the refresh changes \texttt{exposure\_nf\_id}, \texttt{tw\_id}, \texttt{state\_decl\_id}, the accountant family, or the published horizon row, then the old theorem packet is stale rather than merely rewrapped; a successor state-anonymity card is required before that sentence may be reused.\cite{series:abomoinl,series:release,series:planstateequiv}

\paragraph{Minimal reopen ladder.}
Once the live question is no longer ``what does this exact declared state-anonymity packet license?'', the first reopen owner should be named explicitly rather than inferred.
Reopen State~1A for richer time-to-identification conversions or alternate horizon tables; State~2 for DHT committee-contact floors and abort-safe equalization; State~3 for CPPC manifests, monitoring, and public-alert spend; Operational~A or Operational~B for retry schedules, plan pinning, or runtime drift receipts; Eval~1 for notarized evaluation evidence; and Anonymity~C or Release~A for public claim packaging, lineage, or release-facing receipt questions. If the live issue is the exact public/on-request packet that should carry this card onward, reopen the maintained worked example rather than inventing a local packet schema; if the live issue is theorem-root-to-review routing across that already-fixed card, reopen the maintained series-spine bridge rather than rebuilding that route from memory. Claim-sameness and state-surface equivalence rules remain centralized in Synthesis~18 rather than being redefined here.\cite{series:state1addendum,series:state2,series:state3,series:operationalA,series:operationalB,series:evalnotary,series:abomoinl,series:release,series:workedexample,series:seriesspines,series:planstateequiv}

\section{Takeaways}
\begin{itemize}[leftmargin=*]
\item \textbf{State is an anonymity surface:} any key-dependent influence on selection weights, cache contents, or public scores is a witness stream and must be budgeted.
\item \textbf{Small gaps amplify fast:} repetition turns per-epoch gaps $\Delta$ into near-certain identification on the $1/\Delta^2$ scale; refresh horizons and contractive updates are operationally decisive.
\item \textbf{Stability is the right upper-bound language:} DP-style max-divergence bounds on the control-plane kernel export directly to odds-inflation and effective anonymity-set statements used elsewhere in this series.
\item \textbf{Monitoring must be closed-view:} publishable alarms and dashboards are themselves releases; prefer closed-view auditing (AnonDHT State~3) and commit to what is released via receipts (Release~A).
\end{itemize}

\begin{thebibliography}{99}\small

\bibitem{series:schedcompiler}
Anonymous.
\newblock \emph{Scheduling and Compiler Techniques for Metadata-Hiding Overlay Lookups}.
\newblock Series paper (published), 2026.
\newblock Wiki: \texttt{[[2026.01.22 - Mathematics: Scheduling and Compiler Techniques for Metadata-Hiding Overlay Lookups]]}.

\bibitem{series:spectral}
Anonymous.
\newblock \emph{Spectral Anonymity and Optimal Distinguishing Bounds for Random-Walk Delegation in (Possibly Directed) Overlays}.
\newblock Series paper (published), 2026.
\newblock Wiki: \texttt{[[2026.01.23 - Mathematics: Spectral Anonymity and Optimal Distinguishing Bounds for Random-Walk Delegation in (Possibly Directed) Overlays]]}.

\bibitem{series:prefixcapacities}
Anonymous.
\newblock \emph{Prefix Capacities and Separation Cuts for Shared-Kernel Termination-Time Hiding}.
\newblock Series paper (published), 2026.
\newblock Wiki: \texttt{[[2026.01.25 - Mathematics: Prefix Capacities and Separation Cuts for Shared-Kernel Termination-Time Hiding]]}.

\bibitem{series:knapsacktransport}
Anonymous.
\newblock \emph{Optimal Poset-Feasible Padding: Knapsack, Transport, and Dual Certificates for TV Privacy}.
\newblock Series paper (published), 2026.
\newblock Wiki: \texttt{[[2026.01.29 - Mathematics: Optimal Poset-Feasible Padding: Knapsack, Transport, and Dual Certificates for TV Privacy]]}.

\bibitem{series:stoptimeaddendum}
Anonymous.
\newblock \emph{Stop-Time Padding Addendum (unique-only): Max-Mixture Separations and Tail-Sign Deadline Normal Forms}.
\newblock Series note (published), 2026.
\newblock Wiki: \texttt{[[2026.01.26 - Mathematics: Stop-Time Padding Addendum, Max-Mixture Separations and Tail-Sign Deadline Normal Forms]]}.

\bibitem{series:manytestinglowerbounds}
Anonymous.
\newblock \emph{Many-Testing Lower Bounds for Low-Latency Privacy on Hierarchical Overlays}.
\newblock Series paper (published), 2026.

\bibitem{series:state2}
Anonymous.
\newblock Committee Contact-Set Privacy in Anonymous DHT Lookups: MUCC, Availability Floors, and Abort-Safe Equalization.
\newblock Companion preprint in this archive (AnonDHT~State~2).

\bibitem{series:state3}
Anonymous.
\newblock Closed-View Auditing for Stateful Anonymity: Odometers, Alarms, and Machine-Checkable CPPCs.
\newblock Companion preprint in this archive (AnonDHT~State~3).

\bibitem{series:deployedsurfaces}
Anonymous.
\newblock Deployed Observation Surfaces for Anonymous DHTs: Control-plane knobs in IPFS/libp2p delegated routing and OHTTP evolution.
\newblock Series synthesis note (Synthesis~30), 2026. (This archive.)

\bibitem{cao2026nymrep}
Xinmu Alexis Cao and Matthew Green.
\newblock Analysis and Attacks on the Reputation System of {Nym}.
\newblock Proceedings on Privacy Enhancing Technologies (PoPETs), 2026 (Issue~2).
\newblock Cryptology ePrint Archive, Report 2026/101.
\newblock \url{https://eprint.iacr.org/2026/101}
\newblock Author page (links to artifact): \url{https://alexiscao.github.io/}

\bibitem{chaum1981untraceable}
David Chaum.
\newblock Untraceable Electronic Mail, Return Addresses, and Digital Pseudonyms.
\newblock Communications of the ACM, 24(2), 84--90, 1981.

\bibitem{dingledine2004tor}
Roger Dingledine and Nick Mathewson and Paul Syverson.
\newblock Tor: The Second-Generation Onion Router.
\newblock Proceedings of the 13th USENIX Security Symposium, 2004.

\bibitem{piotrowska2017loopix}
Ania M. Piotrowska and Jamie Hayes and Tariq Elahi and Sebastian Meiser and George Danezis.
\newblock The {L}oopix Anonymity System.
\newblock USENIX Security Symposium, 2017.

\bibitem{oldenburg2024mixmatch}
Lennart Oldenburg and Marc Juarez and Enrique Argones R{\'u}a and Claudia Diaz.
\newblock MixMatch: Flow Matching for Mixnet Traffic.
\newblock Proceedings on Privacy Enhancing Technologies, 2024(2), 276--294, 2024.
\newblock DOI: 10.56553/popets-2024-0050.
\newblock \url{https://petsymposium.org/popets/2024/popets-2024-0050.pdf}

\bibitem{dworkroth2014foundations}
Cynthia Dwork and Aaron Roth.
\newblock The Algorithmic Foundations of Differential Privacy.
\newblock Foundations and Trends in Theoretical Computer Science, 9(3--4), 211--407, 2014.

\bibitem{mcsherry2007mechanism}
Frank McSherry and Kunal Talwar.
\newblock Mechanism Design via Differential Privacy.
\newblock 48th Annual {IEEE} Symposium on Foundations of Computer Science (FOCS), 94--103, 2007.
\newblock DOI: 10.1109/FOCS.2007.41.
\newblock Introduces the exponential mechanism viewpoint used widely for DP-style selection.
\newblock \url{https://kunaltalwar.org/papers/expmech.pdf}

\bibitem{dwork2010continual}
Cynthia Dwork and Moni Naor and Toniann Pitassi and Guy N. Rothblum.
\newblock Differential Privacy Under Continual Observation.
\newblock Proceedings of the 42nd {ACM} Symposium on Theory of Computing (STOC), 715--724, 2010.
\newblock DOI: 10.1145/1806689.1806787.
\newblock \url{https://dl.acm.org/doi/10.1145/1806689.1806787}

\bibitem{chan2010continual}
T.-H. Hubert Chan and Elaine Shi and Dawn Song.
\newblock Private and Continual Release of Statistics.
\newblock International Colloquium on Automata, Languages, and Programming (ICALP), Part {II}, 6199, 405--417, 2010.
\newblock Also available as IACR ePrint 2010/076.
\newblock \url{https://people.eecs.berkeley.edu/~dawnsong/papers/2010%20private%20continual.pdf}

\bibitem{rogers2016odometers}
Ryan Rogers and Aaron Roth and Jonathan Ullman and Salil Vadhan.
\newblock Privacy Odometers and Filters: Pay-as-you-Go Composition.
\newblock Advances in Neural Information Processing Systems (NeurIPS), 2016.

\bibitem{lecuuyer2021practical}
Mathias L{\'e}cuyer.
\newblock Practical Privacy Filters and Odometers with R{\'e}nyi Differential Privacy and Applications to Differentially Private Deep Learning.
\newblock 2021.
\newblock arXiv:2103.01379

\bibitem{mironov2017rdp}
Ilya Mironov.
\newblock R{\'e}nyi Differential Privacy.
\newblock 2017 IEEE 30th Computer Security Foundations Symposium (CSF), 2017.
\newblock Also available as arXiv:1702.07476

\bibitem{balle2020rdp_testing}
Borja Balle and Gilles Barthe and Marco Gaboardi and Justin Hsu and Tetsuya Sato.
\newblock Hypothesis Testing Interpretations and R{\'e}nyi Differential Privacy.
\newblock International Conference on Artificial Intelligence and Statistics (AISTATS), 2020.

\bibitem{kairouz2015composition}
Peter Kairouz and Sewoong Oh and Pramod Viswanath.
\newblock The Composition Theorem for Differential Privacy.
\newblock International Conference on Machine Learning (ICML), 2015.
\newblock Also available as arXiv:1311.0776

\bibitem{chan2011continual_observation}
Tianhao Chan and Elaine Shi and Dawn Song.
\newblock Private and Continual Release of Statistics.
\newblock ACM Transactions on Information and System Security, 2011.
\newblock Early conference version in 2010.

\bibitem{wald_sprt_1945}
Abraham Wald.
\newblock Sequential Tests of Statistical Hypotheses.
\newblock Annals of Mathematical Statistics, 16(2), 117--186, 1945.

\bibitem{howard2021confseq}
Steven R. Howard and Aaditya Ramdas and Jon McAuliffe and Jasjeet Sekhon.
\newblock Time-uniform, nonparametric, nonasymptotic confidence sequences.
\newblock Annals of Statistics, 49(2), 1055--1080, 2021.
\newblock DOI: 10.1214/20-AOS1991.
\newblock Also available as arXiv:1810.08240.
\newblock \url{https://projecteuclid.org/journals/annals-of-statistics/volume-49/issue-2/Time-uniform-nonparametric-nonasymptotic-confidence-sequences/10.1214/20-AOS1991.pdf}

\bibitem{ramdas2024evalues}
Aaditya Ramdas and Ruodu Wang.
\newblock Hypothesis Testing with {E}-values.
\newblock Foundations and Trends\textsuperscript{\textregistered} in Statistics, 1(1--2), 1--390, 2025.
\newblock DOI: 10.1561/3600000002.
\newblock Also available as arXiv:2410.23614.

\bibitem{ramdas2023savi}
Aaditya Ramdas and Peter Gr{\"u}nwald and Vladimir Vovk and Glenn Shafer.
\newblock Game-Theoretic Statistics and Safe Anytime-Valid Inference.
\newblock 2023.
\newblock arXiv:2210.01948

\bibitem{Farokhi2020DiscountedDP}
Farhad Farokhi.
\newblock Temporally Discounted Differential Privacy for Evolving Datasets on an Infinite Horizon.
\newblock Proceedings of the 2020 ACM/IEEE 11th International Conference on Cyber-Physical Systems (ICCPS), 2020.
\newblock DOI: 10.1109/ICCPS48487.2020.00008.
\newblock Conference version with DOI; open preprint on arXiv.
\newblock \url{https://arxiv.org/abs/1908.03995}

\bibitem{AnderssonEtAl2024GradualExpiration}
Joel Daniel Andersson and Monika Henzinger and Rasmus Pagh and Teresa Anna Steiner and Jalaj Upadhyay.
\newblock Continual Counting with Gradual Privacy Expiration.
\newblock Advances in Neural Information Processing Systems (NeurIPS), 2024.
\newblock DOI: 10.48550/arXiv.2406.03802.
\newblock NeurIPS 2024; see also NeurIPS proceedings entry (ACM DL: 10.5555/3737916.3738215).
\newblock \url{https://arxiv.org/abs/2406.03802}

\bibitem{xu2024online_evalues}
Ziyu Xu and Aaditya Ramdas.
\newblock Online multiple testing with e-values.
\newblock Proceedings of the 27th International Conference on Artificial Intelligence and Statistics (AISTATS), 238, 2024.
\newblock PMLR 238.

\bibitem{nasr2018deepcorr}
Milad Nasr and Alireza Bahramali and Amir Houmansadr.
\newblock DeepCorr: Strong Flow Correlation Attacks on Tor Using Deep Learning.
\newblock Proceedings of the 2018 ACM SIGSAC Conference on Computer and Communications Security (CCS), 2018.
\newblock DOI: 10.1145/3243734.3243824.
\newblock Also available as arXiv:1808.07285.

\bibitem{oh2022deepcoffea}
Se Eun Oh and Taiji Yang and Nate Mathews and James K. Holland and Mohammad Saidur Rahman and Nicholas Hopper and Matthew Wright.
\newblock DeepCoFFEA: Improved Flow Correlation Attacks on Tor via Metric Learning and Amplification.
\newblock Proceedings of the IEEE Symposium on Security and Privacy (SP), 1915--1932, 2022.
\newblock DOI: 10.1109/SP46214.2022.9833801.

\bibitem{howard2018cher}
Howard, Steven R. and Ramdas, Aaditya and McAuliffe, Jon and Sekhon, Jasjeet.
\newblock Time-uniform Chernoff bounds via nonnegative supermartingales.
\newblock 2018.
\newblock arXiv:1808.03204

\bibitem{wang2020ebh}
Wang, Ruodu and Ramdas, Aaditya.
\newblock False discovery rate control with e-values.
\newblock 2020.
\newblock arXiv:2009.02824 (published in JRSSB 2022)

\bibitem{lopes2024sumo}
Daniela Lopes and Jin-Dong Dong and Pedro Medeiros and Daniel Castro and Diogo Barradas and Bernardo Portela and Jo{\~a}o Vinagre and Bernardo Ferreira and Nicolas Christin and Nuno Santos.
\newblock Flow Correlation Attacks on {Tor} Onion Service Sessions with Sliding Subset Sum.
\newblock Network and Distributed System Security Symposium (NDSS), 2024.
\newblock DOI: 10.14722/ndss.2024.24337.
\newblock \url{https://www.ndss-symposium.org/ndss-paper/flow-correlation-attacks-on-tor-onion-service-sessions-with-sliding-subset-sum/}

\bibitem{wang2022ebh}
Ruodu Wang and Aaditya Ramdas.
\newblock False discovery rate control with e-values.
\newblock Journal of the Royal Statistical Society: Series B, 84(3), 2022.
\newblock doi:10.1111/rssb.12489 (preprint arXiv:2009.02824)

\bibitem{fischer2024onlinebh}
Lasse Fischer and Ziyu Xu and Aaditya Ramdas.
\newblock An online generalization of the (e-)Benjamini--Hochberg procedure.
\newblock 2024.
\newblock arXiv:2407.20683

\bibitem{wang2025sebH}
Wang, Hongjian and Dandapanthula, Sanjit and Ramdas, Aaditya.
\newblock Anytime-valid {FDR} control with the stopped e-{BH} procedure.
\newblock Statistics \& Probability Letters, 226, 110512, 2025.
\newblock DOI: 10.1016/j.spl.2025.110512.
\newblock Preprint: arXiv:2502.08539.

\bibitem{singh2025revealnet}
Gurjot Singh and Alim Dhanani and Diogo Barradas.
\newblock {RevealNet}: Distributed Traffic Correlation for Attack Attribution on Programmable Networks.
\newblock 2025.
\newblock arXiv:2505.00618 (accepted to IEEE NCA 2025)

\bibitem{whitehouse2023fullyadaptive}
Justin Whitehouse and Aaditya Ramdas and Ryan Rogers and Steven Wu.
\newblock Fully-Adaptive Composition in Differential Privacy.
\newblock Proceedings of the 40th International Conference on Machine Learning (ICML), 202, 36990--37007, 2023.
\newblock ICML 2023 / PMLR; an open PDF is also mirrored by the authors.
\newblock \url{https://proceedings.mlr.press/v202/whitehouse23a.html}

\bibitem{zhang2023concurrentinteractive}
Haney, Samuel and Shoemate, Michael and Tian, Grace and Vadhan, Salil and Vyrros, Andrew and Xu, Vicki and Zhang, Wanrong.
\newblock Concurrent Composition for Interactive Differential Privacy with Adaptive Privacy-Loss Parameters.
\newblock Proceedings of the 2023 ACM SIGSAC Conference on Computer and Communications Security (CCS '23), 2023.
\newblock DOI: 10.1145/3576915.3623128.
\newblock Also available as arXiv:2309.05901.

\bibitem{blierwong2026thresholds}
Christopher Blier-Wong and Ruodu Wang.
\newblock Improved thresholds for e-values.
\newblock 2026.
\newblock DOI: 10.48550/arXiv.2408.11307.
\newblock Updated preprint (Feb 7, 2026); see also earlier arXiv versions.
\newblock \url{https://arxiv.org/abs/2408.11307}

\bibitem{vadhanwang2021concurrentdp}
Salil Vadhan and Tianhao Wang.
\newblock Concurrent Composition of Differential Privacy.
\newblock 2021.
\newblock arXiv:2105.14427

\bibitem{lyu2022interactivecomp}
Xuechen Lyu and Zhanqiu Zhang and Ziyu Zhang and Yu-Xiang Wang.
\newblock Composition Theorems for Interactive Differential Privacy.
\newblock Advances in Neural Information Processing Systems (NeurIPS), 2022.
\newblock neurips.cc; 2022

\bibitem{balle2018subsampling}
Borja Balle and Gilles Barthe and Marco Gaboardi.
\newblock Privacy Amplification by Subsampling: Tight Analyses via Couplings and Divergences.
\newblock Advances in Neural Information Processing Systems (NeurIPS), 2018.
\newblock NeurIPS 2018; also available as arXiv:1807.01647

\bibitem{kifer2014pufferfish}
Daniel Kifer and Ashwin Machanavajjhala.
\newblock Pufferfish: A Framework for Mathematical Privacy Definitions.
\newblock ACM Transactions on Database Systems, 39(1), 1--36, 2014.
\newblock DOI: 10.1145/2514689.
\newblock TODS 39(1), Article 3; open author PDF.
\newblock \url{https://users.cs.duke.edu/~ashwin/pubs/pufferfish_TODS.pdf}

\bibitem{Sendner2024MirageFlow}
Christoph Sendner and Jasper Stang and Alexandra Dmitrienko and Raveen Wijewickrama and Murtuza Jadliwala.
\newblock MirageFlow: A New Bandwidth Inflation Attack on {Tor}.
\newblock Network and Distributed System Security Symposium (NDSS), 2024.
\newblock NDSS 2024; PDF also mirrored by anonbib/freehaven.
\newblock \url{https://www.ndss-symposium.org/ndss-paper/mirageflow-a-new-bandwidth-inflation-attack-on-tor/}

\bibitem{Sendner2023TorMult}
Christoph Sendner and Jasper Stang and Alexandra Dmitrienko and Raveen Wijewickrama and Murtuza Jadliwala.
\newblock TorMult: Introducing a Novel {Tor} Bandwidth Inflation Attack.
\newblock 2023.
\newblock DOI: 10.48550/arXiv.2307.08550.
\newblock arXiv:2307.08550
\newblock \url{https://arxiv.org/abs/2307.08550}

\bibitem{Hanley2019DPSelect}
Hans Hanley and Yixin Sun and Sameer Wagh and Prateek Mittal.
\newblock DPSelect: A Differential Privacy Based Guard Relay Selection Algorithm for Tor.
\newblock Proceedings on Privacy Enhancing Technologies, 2019(2), 166--186, 2019.
\newblock DOI: 10.2478/popets-2019-0025.
\newblock PoPETs 2019 (open access via PETS).
\newblock \url{https://petsymposium.org/popets/2019/popets-2019-0025.php}

\bibitem{pets2026_paperlist_nymrep}
{Privacy Enhancing Technologies Symposium (PETS)}.
\newblock PETS 2026 Accepted Papers List.
\newblock 2026.
\newblock Accepted paper list; includes ``Analysis and Attacks on the Reputation System of Nym'' by Xinmu Alexis Cao and Matthew Green.
\newblock \url{https://petsymposium.org/2026/paperlist.php}

\bibitem{jolles2025wf4nym_zenodo}
Eric Joll{\`e}s.
\newblock Website Fingerprinting on {Nym}: Attacks and Defenses, Artifact Test Data.
\newblock Zenodo dataset, 2025.
\newblock DOI: 10.5281/zenodo.17867461.
\newblock Published 2025-12-09; artifact dataset for the PETS 2026 paper ``Website Fingerprinting on Nym: Attacks and Defenses''.
\newblock \url{https://zenodo.org/records/17867461}

\bibitem{nym_cryptoecon_reward_sharing}
Diaz, Claudia and Halpin, Harry and Kiayias, Aggelos.
\newblock Reward Sharing for Mixnets.
\newblock 2022.
\newblock Nym cryptoeconomic paper; notes ``720 hourly epochs per monthly interval'' while cautioning epoch/interval lengths may change.
\newblock \url{https://nym.com/nym-cryptoecon-paper.pdf}


\bibitem{series:artifactpolicy}
Anonymous.
\newblock Artifact and Disclosure Policy for the One-True-Archive (Source-Only Public Release).
\newblock Synthesis~6, The-One-True-Archive (2026).

\bibitem{series:notation}
Anonymous.
\newblock Notation and Minimal Model for the One-True-Archive.
\newblock Synthesis~8, The-One-True-Archive (2026).

\bibitem{series:contractobjects}
Anonymous.
\newblock Contract Objects for Anonymous-DHT Privacy Budgets and Claims.
\newblock Synthesis~0, The-One-True-Archive (2026).

\bibitem{series:abomoinl}
Anonymous.
\newblock ABOM and OINL for Anonymity Claims.
\newblock Anonymity series Paper~C, The-One-True-Archive (2026).

\bibitem{series:oddsinflation}
Anonymous.
\newblock Odds Inflation, Privacy Loss, and Endpoint Anonymity Bounds.
\newblock Anonymity series Paper~A, The-One-True-Archive (2026).

\bibitem{series:endpointbridge}
Anonymous.
\newblock Endpoint Metrics Bridge: Odds Inflation, PML/MaxL, and Identification Sizing.
\newblock Synthesis~5, The-One-True-Archive (2026).

\bibitem{series:traceifaces}
Anonymous.
\newblock Trace Interfaces for Anonymous DHTs.
\newblock Synthesis~3, The-One-True-Archive (2026).

\bibitem{series:threatwindows}
Anonymous.
\newblock Threat Models and Repetition Windows for Anonymous DHTs.
\newblock Synthesis~4, The-One-True-Archive (2026).

\bibitem{series:receiptitems}
Anonymous.
\newblock Receipt Line Items for Anonymity Claims: A Minimal Tuple Schema for Published Budgets.
\newblock Synthesis~12, The-One-True-Archive (2026).

\bibitem{series:planstateequiv}
Anonymous.
\newblock Digest-Bound Replay Plans and State Declarations: Equivalence, Minimality, and Change Control.
\newblock Synthesis~18, The-One-True-Archive (2026).

\bibitem{series:workedexample}
Anonymous.
\newblock Worked Example: Instantiating a Tiered Reader-Privacy Receipt.
\newblock Synthesis~17, The-One-True-Archive (2026).

\bibitem{series:seriesspines}
Anonymous.
\newblock Series Spines from Mathematics to Review: Stable Handoffs for Later Terse Papers.
\newblock Synthesis~55, The-One-True-Archive (2026).

\bibitem{series:evalnotary}
Anonymous.
\newblock Notarized Anonymity Certificates: Anti-Grinding Receipts for Evaluation Attempts.
\newblock Evaluation series Paper~1, The-One-True-Archive (2026).

\bibitem{series:operationalA}
Anonymous.
\newblock Schedule-Only Retries in Anonymous DHTs: When Retry Logic Becomes a Linkability Channel.
\newblock Operational series Paper~A, The-One-True-Archive (2026).

\bibitem{series:operationalB}
Anonymous.
\newblock Auditable Trace Privacy with Abort: Plan-Pinning and Drift Receipts for Side-Channel Resistant Transport.
\newblock Operational series Paper~B, The-One-True-Archive (2026).

\bibitem{series:release}
Anonymous.
\newblock Release Property, Ledger Receipts, and Public Lineage for Anonymity Claims.
\newblock Release-and-Destination series Paper~A, The-One-True-Archive (2026).

\bibitem{series:state1addendum}
Anonymous.
\newblock State-Dependent Anonymity (Addendum): Calibration Vignettes and Time-to-Identification Conversions.
\newblock AnonDHT State series Paper~1A, The-One-True-Archive (2026).


\bibitem{series:condcert}
Anonymous.
\newblock Conditional Per-Step Budget Certificates.
\newblock Series synthesis note (Synthesis~21), 2026. (This archive.)



% Added in rev0806 citation-closure hygiene pass.
\bibitem{series:state3addendum}
Anonymous.
\newblock Closed-View Auditing and CPPCs (Addendum): Multi-Stream Monitoring Notes and Schema Excerpts.
\newblock AnonDHT State series Paper~3A, 2026. (This archive.)

\end{thebibliography}

\end{document}
